% Lösungen Unabhängigkeit (zusammenbau v0.4, Heft)
\einheitenkopf{Unabhängigkeit}

\erg{21}{a) alle drei stehen in der Tafel: zwei Ränder und ein Feld – $P(A)$ und $P(B)$ am Rand, $P(A \cap B)$ im Feld \quad b) $P(M \cap F) = \frac{20}{100} = 0{,}2$, $P(M) \cdot P(F) = 0{,}4 \cdot 0{,}5 = 0{,}2$ – gleich, also unabhängig \quad c) 12 Erwachsene spielen Volleyball; $P(E \cap V) = \frac{12}{60} = 0{,}2$, $P(E) \cdot P(V) = 0{,}6 \cdot \frac{1}{3} = 0{,}2$ – gleich, also unabhängig}

21c)

\vierfeldertafel{E}{V}{12,8,20,24,16,40,36,24,60}

\erg{22}{a) $P(S \cap H) = \frac{70}{250} = 0{,}28$, $P(S) \cdot P(H) = 0{,}4 \cdot 0{,}7 = 0{,}28$ – gleich, also unabhängig \quad b) Ältere mit Bestehen: $9\,920 - 7\,750 = 2\,170$; $P_{Ä}(S) = \frac{2\,170}{3\,100} = 0{,}7$, $P(S) = \frac{9\,920}{12\,400} = 0{,}8$ – ungleich, also abhängig: Ältere bestehen seltener als der Durchschnitt. \quad c) schwer: $\frac{90}{150} = 0{,}6$; nicht schwer: $\frac{150}{450} \approx 0{,}333$ – verschieden, also abhängig}
\erg{23}{a) $0{,}5$ – die Münze hat kein Gedächtnis, jeder Wurf ist unabhängig \quad b) Falsch: Jede Drehung ist unabhängig von den vorigen, die Wahrscheinlichkeit für Grün bleibt die Hälfte; das Rad muss nichts ausgleichen.}
\erg{24}{a) Falsch: Die Samstage sind voneinander unabhängig, die Wahrscheinlichkeit für sieben richtige Antworten ist an jedem Samstag dieselbe; die früheren Ergebnisse ändern sie nicht – dass vier solche Samstage nacheinander selten sind, ist kein Argument für den einzelnen Samstag. \quad b) Falsch: Die Serien sind voneinander unabhängig, die Wahrscheinlichkeit für mindestens einen gehaltenen Elfmeter ist in jeder Serie gleich; die drei Serien ohne Erfolg erhöhen sie nicht.}
